\subsection{\texttt{TNOT}}
\noindent\textbf{Bundle encoding.} UOP entry 33, opcode \texttt{0x103E}. Bundle: \texttt{B.OP + B.ARG + B.IOTI}. \texttt{B.OP.Opcode}=\texttt{0x103E}. \texttt{B.IOTI}: selector=\texttt{111 last}, \texttt{SrcTile0}=\texttt{src}, \texttt{SrcTile1}=\texttt{absent}, \texttt{DstTile}=\texttt{dst}, \texttt{SizeCode}=profile tile size.\par\smallskip
{\small
\paragraph{PTO ISA description.}
Elementwise bitwise NOT of a tile.
\paragraph{Example.}
i8 tile [0b00001010, 0b11110000] gives [0b11110101, 0b00001111].
\par\smallskip
\paragraph{PTO ISA operation.}
For each element \texttt{(i, j)} in the valid region:

\[
\mathrm{dst}_{i,j} = \sim\mathrm{src}_{i,j}
\]
}\par\smallskip
\paragraph{Notes.}
For FP types, the result is the bitwise NOT applied to the IEEE 754 encoding. This is NOT equivalent to multiplying by -1.
\paragraph{Microcode site table.}
{\scriptsize
\begin{longtable}{@{}R{0.055\linewidth}L{0.23\linewidth}L{0.31\linewidth}L{0.22\linewidth}@{}}
\toprule
Seq & Micro instruction & WO[63:32] fields & WM[31:0] fields \\
\midrule
0 & \texttt{u\_vec.vlds} & \texttt{opc=0x17, x=0, fl=mem, seq=0, kind=b} & \texttt{entry=33, cnt=2, res=vec, var=0} \\
1 & \texttt{u\_vec.vnot} & \texttt{opc=0x21, x=0, fl=alu, seq=1, kind=u} & \texttt{entry=33, cnt=1, res=vec, var=0} \\
2 & \texttt{u\_vec.vsts} & \texttt{opc=0x2A, x=0, fl=mem, seq=2, kind=b} & \texttt{entry=33, cnt=2, res=vec, var=0} \\
\bottomrule
\end{longtable}
}
\paragraph{Microcode body DSL.}
\begin{lstlisting}[style=ptocode]
def uop_tnot(src, dst):
    dtype = dst.element_type
    valid_rows, valid_cols = dst.valid_shape

    for row in range(0, valid_rows, 1):
        remained = valid_cols
        for col in range(0, valid_cols, pto.get_lanes(dtype)):
            mask, remained = pto.make_mask(dtype, remained)
            vinput = pto.vlds(src[row, col:])
            result = pto.vnot(vinput, mask)
            pto.vsts(result, dst[row, col:], mask)
    return
\end{lstlisting}
\Needspace{0.34\textheight}
